% ============================================================================
% ICSE Paper — Approach Section (v5: updated for S-Linker12e)
%
% Changes from v4:
%   - Layer 1 now runs 3 analyses in parallel (added seed extraction);
%     no 1a/1b sub-tier split
%   - Trailing-word enrichment merged into document analysis
%   - Seed validation replaces evidence-stratified revalidation:
%     per-component disambiguation (COMPONENT/OTHER) with component
%     profile, anchor sentences, and approve-biased framing
%   - Entity validation uses pure intersection (no alias stratification)
%   - Abbreviated reference agent removed (trailing words feed into
%     entity extraction via alias table)
%   - Consolidation simplified to priority-ordered dedup (3 agents)
%   - Table and figure updated for 3-agent Layer 2
% ============================================================================

\section{AALinker}\label{sec:approach}

We propose a multi-agent workflow for recovering trace links between
architectural documents and model components.
Given an architecture document~$D$ consisting of sentences
$S = \{s_1, \ldots, s_n\}$ and an architectural model containing
components $C = \{c_1, \ldots, c_m\}$, our approach produces a set of
trace links $L \subseteq S \times C$, where each link $(s_i, c_j)$
asserts that sentence~$s_i$ discusses component~$c_j$'s architectural
responsibilities.

Two problems make this task difficult.
First, \emph{reference diversity}: documents refer to components by
canonical name, abbreviation, functional description, trailing word,
or pronoun.
No single extraction strategy handles all of these forms.
Second, \emph{name ambiguity}: component names like \texttt{Cache} or
\texttt{Facade} are common English words that appear frequently
without referring to any component, so broadening extraction to
capture diverse references also increases false positives.

Addressing both problems requires structured knowledge: which names
are ambiguous, and what abbreviations, synonyms, and short forms map
to which component.
To the best of our knowledge, this knowledge does not exist in advance
and must be discovered at runtime, using the same LLM that will
consume it, with the risk of hallucinated mappings or misjudged
ambiguity.
We call this the \emph{architectural context knowledge gap}.

Existing approaches address only part of this problem.
Lexical NLP pipelines such as SWATTR~\cite{swattr} rely on string
similarity and cannot resolve indirect references or discourse
dependencies.
LLM-based classifiers such as LISSA~\cite{lissa} add semantic
understanding but have no mechanism to build alias mappings or classify
name ambiguity before linking, treat each sentence in isolation, and
do not address LLM non-determinism.

Our pipeline closes the knowledge gap first, then leverages multiple
link-recovery agents that consume the acquired knowledge.
It is organized as a DAG (Figure~\ref{fig:pipeline}) with three
layers:

\begin{enumerate}
\item \textbf{Project knowledge discovery} (Layer~1):
  Three analyses run in parallel: a \emph{model understanding} agent
  classifies component name ambiguity, a \emph{document understanding}
  agent discovers component aliases---including trailing-word
  references---and a \emph{seed extraction} agent recovers an initial
  set of explicit-reference links via two complementary LLM passes
  (\S\ref{sec:knowledge}).
  Seed extraction requires no Layer~1 knowledge, so it runs
  concurrently with the knowledge analyses.

\item \textbf{Candidate link recovery} (Layer~2):
  Three complementary agents refine and extend the initial link set:
  \emph{seed reference disambiguation} filters false positives from
  the seed using the acquired knowledge,
  \emph{contextual entity extraction} recovers alias-based and
  functional references,
  and \emph{anaphoric coreference resolution} resolves pronominal
  references.
  The agents run in parallel, each targeting a different type of
  component reference
  (Table~\ref{tab:strategies}, \S\ref{sec:recovery}).

\item \textbf{Link consolidation} (Layer~3):
  All recovered links are merged into the final set by deduplicating
  overlaps with first-seen priority (\S\ref{sec:merge}).
\end{enumerate}

Since our agents are LLM-based, their outputs may contain noise such
as hallucinated results and non-deterministic variation across runs.
To handle them, our pipeline uses two cross-cutting
patterns.~\todo{cite source}
\emph{Multi-pass consensus}: for open-ended tasks, we run two
independent LLM passes with task-specific agreement rules to filter
single-run variance.
\emph{Citation-grounded verification}: for verifiable tasks, the LLM
must produce checkable evidence---such as antecedent sentence
numbers---that deterministic code can verify.

\begin{table}[t]
\caption{Link recovery agents and their properties.}
\label{tab:strategies}
\centering\footnotesize
\begin{tabular}{@{}lllll@{}}
\toprule
\textbf{Agent} & \textbf{Reference type} & \textbf{Knowledge} & \textbf{Consensus pattern} & \textbf{LLM} \\
\midrule
Explicit & Canonical name & Alias table, ambig.\rlap{$^\ddagger$} & Asym.\ voting + disambig. & 2+1/comp\rlap{$^*$} \\
Contextual & Alias, functional & Alias table, ambig.\ classif. & $\cap$ extr.\ + $\cap$ valid. & 2+2--3\rlap{$^\dagger$} \\
Anaphoric & Pronoun & Alias table & Citation verif. & 1 pass \\
\bottomrule
\end{tabular}
\\[2pt]
{\scriptsize $^*$2 extraction passes (Layer~1) + 1~disambiguation pass per component (Layer~2).
$^\dagger$2 extraction + 2--3 validation (generic-mention detection for ambiguous names, then 2 voting passes).
$^\ddagger$Extraction needs no knowledge (runs parallel with Layer~1);
disambiguation uses the alias table and ambiguity classification once Layer~1 completes.}
\end{table}

% ---------------------------------------------------------------------------
% ============================================================================
% TikZ figure for the pipeline DAG
% Requires: \usepackage{tikz}, \usetikzlibrary{positioning,arrows.meta,fit,backgrounds}
% ============================================================================

\begin{figure*}[t]
\centering
\begin{tikzpicture}[
    node distance=0.8cm and 0.6cm,
    phase/.style={draw, rounded corners, minimum width=2.2cm, minimum height=0.7cm,
                  font=\footnotesize, align=center},
    layer1/.style={phase, fill=blue!8},
    layer2/.style={phase, fill=orange!10},
    layer3/.style={phase, fill=green!8},
    arr/.style={-{Stealth[length=2mm]}, thick, gray!70},
    layerbox/.style={draw=gray!40, dashed, rounded corners=3pt, inner sep=6pt},
    layerlabel/.style={font=\scriptsize\sffamily, gray!60},
]

% Layer 1a: Independent analyses (top row)
\node[layer1] (model)   {Model\\Analysis};
\node[layer1, right=1.2cm of model] (doc) {Document\\Analysis};

% Layer 1b: Enrichment (needs 1a results)
\node[layer1, below=0.7cm of model, xshift=1.7cm] (enrich) {Partial-Ref.\\Refinement};

% Layer 1 internal dependencies
\draw[arr] (model.south east) -- (enrich.north west);
\draw[arr] (doc.south west) -- (enrich.north east);

% Layer 2: Link Recovery
\node[layer2, below=1.4cm of enrich, xshift=-2.0cm] (explicit) {Explicit\\Ref.\ Extr.};
\node[layer2, right=0.6cm of explicit] (contextual) {Contextual\\Ref.\ Disc.};
\node[layer2, right=0.6cm of contextual] (anaphoric) {Anaphoric\\Ref.\ Res.};
\node[layer2, right=0.6cm of anaphoric] (abbreviated) {Abbreviated\\Ref.\ Match};

% Knowledge flow arrows (Layer 1 -> Layer 2)
\draw[arr] (enrich.south west) -- (explicit.north east);
\draw[arr] (model.south) -- (contextual.north west);
\draw[arr] (enrich.south) -- (contextual.north);
\draw[arr] (enrich.south east) -- (anaphoric.north);
\draw[arr] (enrich.south east) -- (abbreviated.north west);

% Layer 3: Link Consolidation
\node[layer3, below=1.4cm of contextual, xshift=0.8cm] (merge) {Merge};

% Dedup dependency: abbreviated matching needs prior strategies' results
\draw[arr, dashed] (explicit.east) -- (abbreviated.west);

% Layer 2 -> Layer 3
\draw[arr] (explicit.south) -- (merge.north west);
\draw[arr] (contextual.south) -- (merge.north);
\draw[arr] (anaphoric.south) -- (merge.north);
\draw[arr] (abbreviated.south) -- (merge.north east);

% Output
\node[below=0.4cm of merge, font=\footnotesize\itshape] (out) {Trace Links $L$};
\draw[arr] (merge) -- (out);

% Layer labels
\node[layerlabel, left=0.3cm of model] {Layer 1a};
\node[layerlabel, left=0.3cm of enrich] {Layer 1b};
\node[layerlabel, left=0.3cm of explicit] {Layer 2};
\node[layerlabel, left=0.8cm of merge] {Layer 3};

% Layer boxes
\begin{scope}[on background layer]
    \node[layerbox, fit=(model)(doc)(enrich), label={[layerlabel]above:Knowledge Acquisition}] {};
    \node[layerbox, fit=(explicit)(abbreviated), label={[layerlabel]above left:Link Recovery}] {};
    \node[layerbox, fit=(merge), label={[layerlabel]above left:Link Consolidation}] {};
\end{scope}

\end{tikzpicture}
\caption{Pipeline architecture as a staged DAG. Layer~1 acquires knowledge
in two sub-tiers: model analysis and document analysis run
concurrently~(1a); partial-reference refinement depends on their
results~(1b).
Explicit reference \emph{extraction} has no knowledge dependencies
and starts concurrently with Layer~1 (shown here under Link Recovery
for grouping); its \emph{validation} uses the alias table once
Layer~1 completes.
Layer~2 runs four recovery strategies; abbreviated matching
depends on prior strategies for deduplication (dashed arrow).
Layer~3 deduplicates their output.}
\label{fig:pipeline}
\end{figure*}

% ---------------------------------------------------------------------------

\subsection{Project Knowledge Discovery}\label{sec:knowledge}

Architectural documents contain diverse references that often confuse
link-recovery agents.
In addition, component names are often common English words
(\texttt{Core}, \texttt{Adapter}) or technologies (\texttt{Redis}).
These words appear frequently without referring to the component that
shares the same name, posing a challenge for deep understanding of
both document and model in the architectural context.

Layer~1 builds a shared knowledge base and a first set of candidate
links.
Three analyses run in parallel: a \emph{model understanding agent}
analyzes the architectural model to identify ambiguous components
(\S\ref{sec:model-understanding}), a \emph{document understanding
agent} extracts component aliases from the document---including
LLM-classified trailing-word mappings
(\S\ref{sec:doc-understanding})---and a \emph{seed extraction agent}
recovers explicit-reference links without requiring any knowledge
(\S\ref{sec:seed}).
Once all three complete, the resulting knowledge base---a list of
ambiguous components and an alias mapping table---together with the
raw seed links, is consumed by all Layer~2 agents.

\subsubsection{Architectural Model Understanding}\label{sec:model-understanding}

Architectural components operate at a high level of abstraction, so
their names vary widely in specificity.
Constructed identifiers like \texttt{TaskScheduler} unambiguously name
architectural roles, but single common-English words like
\texttt{Core} or \texttt{Adapter} appear frequently both when
referring to a specific component and when used as general
concepts.~\todo{cite source}
This poses challenges for trace link recovery: blindly matching
mentions of such names would produce false positives, and an LLM
lacking explicit reasoning about name ambiguity may conflate general
usage with component references.
Therefore, we analyze component names from the input architecture model
to identify \emph{ambiguous} components---those whose names can also
be used as general concepts (e.g., \emph{logic}, \emph{storage}).
This classification serves as knowledge to guide downstream agents
when evaluating candidate links.

Motivated by prior work on CamelCase identifiers and
acronyms,~\todo{cite} we combine structural rules with a few-shot
LLM prompt to mark ambiguous components.
Structural cues---CamelCase boundaries, multiple words, all-uppercase
letters---resolve most cases without LLM involvement: names exhibiting
these cues are classified as \emph{architectural} (non-ambiguous).
The pipeline submits remaining single-word names to the LLM with
few-shot examples for \emph{ambiguous} classification.
The result feeds downstream decisions:
which candidates require generic-mention detection
(\S\ref{sec:entity}) and
which trailing words need capitalized matching
(\S\ref{sec:doc-understanding}).

Additionally, the analysis extracts \emph{partial words} shared among
components.
Because architectural models are hierarchical, a document may refer to
a parent component using only the shared trailing word, creating two
sources of ambiguity:
\begin{enumerate}
  \item Two components sharing a trailing word
    (e.g., \texttt{BasicScheduler}~$\to$~\texttt{Scheduler}): when the
    document says only ``Scheduler,'' it is unclear which component is
    meant.
  \item A component and a general concept: the trailing word may be
    generic English rather than a component reference.
    \todo{concrete example}
\end{enumerate}
Such information is used during alias enrichment
(\S\ref{sec:doc-understanding}) and made available to Layer~2 agents.

\subsubsection{Architectural Document Understanding}\label{sec:doc-understanding}

Architectural documents contain heavy use of synonyms and
abbreviations.~\todo{cite support}
This step finds all alternative names for architectural components
so that downstream agents can recognize references that do not use the
canonical component name.

\paragraph{Alternative name discovery.}
An LLM extracts three categories of alternative names:
\begin{itemize}
    \item \emph{Abbreviations}: short forms introduced via parenthetical
      definitions (e.g., ``Abstract Syntax Tree (AST)'' defines AST).
    \item \emph{Synonyms}: alternative names that unambiguously identify
      exactly one component.
      \todo{example from prompt}
    \item \emph{Partial references}: trailing words of multi-word
      component names used alone (e.g., ``Dispatcher'' for
      \texttt{EventDispatcher}).
\end{itemize}

Different alias types offer different validation opportunities.
Abbreviations carry structural evidence---their letters match the
component name's initials---enabling rule-based auto-approval.
Synonyms lack such structure and need semantic judgment, but
CamelCase identifiers can still be auto-approved since they are
constructed names unlikely to be ordinary English.
A calibrated few-shot judge validates the discovered mappings:
the judge is instructed to auto-approve abbreviations whose letters
match the component name's initials and CamelCase identifiers;
remaining candidates are judged semantically, rejecting generic
terms (\emph{system}, \emph{process}).

\paragraph{Trailing-word enrichment.}
Among the three alias types, partial references require the most
careful handling: trailing-word shorthands are pervasive in
architecture documents, but many trailing words are generic English
words whose referential status depends on context.
After the core discovery and judge steps, a trailing-word enrichment
step identifies additional partial-reference aliases that the initial
extraction may have missed.

The step structurally identifies candidates: trailing words of
multi-word component names that (a)~appear in at least one sentence
where the full component name is absent, (b)~are not shared by
multiple components, and (c)~are at least four characters long.
An LLM-based check classifies each candidate by examining the
sentences where the trailing word appears without the full name,
determining whether the word is used as a standalone reference to
the component or as an ordinary English word.
Only candidates classified as component references are added to the
alias table.
The resulting alias table is consumed by all three recovery agents.

\subsubsection{Seed Extraction}\label{sec:seed-extraction}

The simplest form of component reference is the canonical name itself.
A self-contained extractor uses two complementary LLM passes over
batched document segments:
Pass~A (mention-framed) asks which components are explicitly
mentioned; Pass~B (actor-framed) asks which component is the
primary actor.
Both passes reject names inside hierarchical namespace paths and
generic English words.
Exact-name matches from \emph{either} pass are accepted (union);
non-exact matches require agreement from \emph{both}
(intersection)---a rule we call \emph{asymmetric voting}.
This reflects that exact matches are high-confidence and reliably
found by a single pass, while non-exact matches risk hallucination
and benefit from independent agreement.

Because seed extraction requires no Layer~1 knowledge, it runs in
parallel with knowledge acquisition.
The raw seed links are then refined by the seed disambiguation agent
in Layer~2 (\S\ref{sec:seed-disambig}).

% ---------------------------------------------------------------------------

\subsection{Candidate Link Recovery}\label{sec:recovery}

With ambiguity classifications and the enriched alias table in hand,
three agents target distinct reference types using this shared
knowledge base.
All three agents run in parallel.

\subsubsection{Seed Reference Disambiguation}\label{sec:seed-disambig}

The seed extractor produces high-precision links but operates without
knowledge of which names are ambiguous or what aliases exist.
With Layer~1 knowledge now available, the seed links can be refined:
some candidates may refer to code-level identifiers, generic English
words, or techniques that share a component's name rather than the
component itself.

The agent groups raw seed links by component and constructs a
per-component prompt containing three elements:
\begin{enumerate}
  \item A \emph{component profile} summarizing the component's
    ambiguity classification and known aliases from Layer~1.
  \item \emph{Anchor sentences}---up to five sentences elsewhere in the
    document where the component name appears in proper case but is
    \emph{not} in the seed set---providing calibration examples of
    confirmed references.
  \item Each seed link with its sentence text, preceding sentence, and
    a \emph{mention classification} (proper-case standalone,
    lowercase, inside dotted path, via alias, or indirect).
\end{enumerate}

For each seed link, the LLM classifies the reference as
\textsc{component} (the sentence discusses the architectural
component) or \textsc{other} (the name carries a different meaning:
code-level notation, technique, embedded sub-entity, different entity,
or generic English).
The prompt is approve-biased---when uncertain, the LLM is instructed
to choose \textsc{component}---reflecting the high precision of the
seed extractor's asymmetric voting.
Rejected links are removed; approved links proceed to consolidation.

\subsubsection{Contextual Reference Agent}\label{sec:entity}

Alias-based and functional references require knowledge that the
seed extractor lacks, but broadening the extraction scope also
increases the risk of hallucinated links.
Running the \emph{same} broad-scope prompt twice independently
and keeping only shared results filters out hallucinations while
preserving the recall benefit of broader scope.
Using two different prompts would introduce prompt-specific biases;
using the same prompt ensures that disagreements between the two
runs reflect only LLM non-determinism.

Knowledge-enriched prompts receive the known aliases from the
alias table, enabling recognition of abbreviations, synonyms, and
partial names that the seed extractor cannot match.
The extraction scope includes space-separated compound forms,
functional descriptions by name or role, interaction participants
(``X sends data to Y'' covers both X and Y),
passive and prepositional mentions (``handled by X'', ``via X'',
``through X''), and known alias occurrences.
The same prompt runs twice independently, and the pipeline keeps
only candidates found in both passes.

\paragraph{Validation.}
Candidates undergo two-stage validation:
\begin{enumerate}
    \item \emph{Generic-mention detection:}
      For components whose names were classified as ambiguous by
      the model understanding agent (\S\ref{sec:model-understanding}),
      an LLM determines whether each lowercase occurrence is a
      component reference or ordinary English usage, given the
      surrounding sentence context and anchor sentences where the
      full name appears;
      the pipeline rejects ordinary-usage occurrences.
    \item \emph{Two-pass intersection:}
      Two independent LLM passes with complementary
      foci---\emph{architectural participation} (does the sentence
      name this component as an architectural participant performing
      operations or providing services?)
      and \emph{referential specificity} (is the component name used
      to identify this specific architectural element, or does it
      serve as a generic technical term?)---validate remaining
      candidates.
      Each candidate receives an evidence context including matched
      text span, mention type, preceding sentence, and anchor
      sentences from elsewhere in the document.
      A candidate is accepted only if \emph{both} passes approve
      (intersection).
\end{enumerate}

\subsubsection{Anaphoric Reference Agent}\label{sec:coref}

Sentences containing pronouns (\emph{it}, \emph{they}, \emph{this},
\emph{these}, \emph{that}, \emph{those}, \emph{its}, \emph{their})
carry valid trace information but contain no component name, so
the preceding agents cannot find them.
However, LLMs can fabricate antecedents---claiming a pronoun refers
to a component that was never mentioned in nearby sentences.
If the LLM is required to cite a specific antecedent sentence,
such fabrications become detectable: the cited sentence either
contains the component name (or a known alias) or it does not.
A separate validation LLM call would itself be subject to
fabrication; deterministic string verification eliminates this risk.

Each pronoun-bearing sentence is presented with a $\pm5$-sentence
context window.
The broader window provides surrounding text for disambiguation,
while the prompt instructs the LLM to resolve only within the
previous 3~sentences---a conservative bound that reduces
long-distance fabrications.
The LLM must cite a specific antecedent sentence number.
The pipeline verifies this citation by checking that the
component name (or a known alias) appears verbatim in the
cited antecedent sentence.
Because this check is strict and deterministic, anaphoric
links that pass are accepted without further LLM review.

% ---------------------------------------------------------------------------

\subsection{Link Consolidation}\label{sec:merge}

The three recovery agents produce independent candidate
sets, each already validated by its own consensus pattern
(Table~\ref{tab:strategies}).
Because agents run independently, they may discover the
same $(s_i, c_j)$ pair through different evidence.
Overlaps are removed by first-seen priority:
candidates are processed in agent order
(explicit, contextual, anaphoric),
and the first source to discover a pair is kept.
